\scanpage{128}
% archive sheet number: 127
\unsure{+)|*)} $(S,A)$ graphe \unsure{trilitaire|trilatéraire|triangulaire} $\mathcal T_{\nu}$, avec $\deleted{\unsure{\nu>0\in\mathbb N|\nu\in\mathbb N}}\inserted{\nu=0,1\text{ ou }2}$, $c_\nu=2^\nu\geq1$

\[a\in S\]

$I$ l'ens des triangles de sommet $a$ (de cardinal $\deleted{\unsure{e}}c_\nu\inserted{+1}=2^\nu+1$)

$V$ \deleted{le groupe engendré par les $\sigma_{\underline t}$ $(\underline t\in I)$} $\deleted{\unsure{(\mathbb Z/2)^I}}\simeq\mathbb F_2^I$, il opère sur $S$ grâce aux $\sigma_{\underline t}$ $(\underline t\in I)$

$\deleted{\unsure{S|S_A}}X=\{s\in S\mid s\ \text{lié à }a\}$

\deleted{\unsure{$s\longmapsto s$|$s\longmapsto s'$}} Pour $s\in X$, on désigne par $s'$ le \deleted{\unsure{dernier}}\inserted{troisième} sommet du triangle eng. par l'arête $\{a,s\}$

Donc $s\mapsto s'$ est une involution sans pt fixe $\sigma$ de $X$, et $I\simeq X/\sigma$. On peut regarder les $X_i$ $(i\in I)$ \deleted{\unsure{soit}} comme des torseurs sous $\mathbb F_2$, donc $\prod_{i\in I}X_i$ comme un torseur sous $V$, soit $E$

\deleted{Soit} $\deleted{\unsure{E}}Y$ l'ens $S\smallsetminus X\smallsetminus\{a\}=\{\deleted{\unsure{x\in S}}\inserted{x\in S\smallsetminus\{a\}}\mid x\ \text{non lié à }a\inserted{\text{ et }a\ne}\}$

Pour tt $x\in Y$ et $i\in I$, $x$ est lié à un elt de $E_i$ et un seul, soit $\varphi_i(x)$, d'où
\[
\varphi(x)=(\varphi_i(x))_{i\in I}\in E=\prod X_i.
\]
d'où \deleted{Comment $G$ opère-t}
\[
\varphi:\ Y\longrightarrow E
\]

Comment $G$ opère-t-il (via $G\to\operatorname{Aut}(S)$.) sur $E$ ? Soit $\deleted{z\text{''}}$ \inserted{$z_0$ l'élt non nul de $\mathbb F_2$, $z$ son image \unsure{ds|de} $\mathbb F_2^I$} \deleted{l'élémt diagonal \illegible{1} de $\mathbb F_2^I$} \inserted{via l'application diagonale.} \deleted{correspond à $\unsure{t\in I}\quad\unsure{t}$ Alors $\sigma_i$ $(i\in I)$}

soit $\varepsilon_i:\mathbb F_2\longrightarrow V\simeq\mathbb F_2^I$ \deleted{\unsure{défini}} l'inj. canonique, $\rho$ l'op. naturelle \deleted{de} $V$ sur $E$ par $\rho\varepsilon_i(z_0):(x_j)_j=(y_j)_j$ avec $y_j=x_j$ si $j\ne i$, $y_j=\unsure{\varepsilon_0|z_0}x_j$ sinon. On aura alors $\deleted{\sigma_t(x)=}$
